\section{System Design}
A \emph{precompile} is a dedicated sub-circuit (an \emph{AIR}: algebraic intermediate representation, the STARK trace's constraint form) for a designated primitive that the zkVM guest invokes as a syscall, so the primitive's cost bypasses field emulation. We construct one around a reusable interface and instantiate it as a small \emph{family} of three chips over PLUM's $199$-bit prime field: (i) a bespoke Griffin permutation~\cite{cryptoeprint:2022/403} AIR, the load-bearing $91\%$ target; (ii) a dedicated \texttt{FP192\_MUL} field-multiplication chip; and (iii) a \texttt{FP192\_POW\_RES} chip composing that multiplication into the $t$-th power-residue PRF symbol check $a^{(p-1)/t}\bmod p$ ($t=256$ fixed). All three build and pass a core-STARK smoke proof (execute, prove, verify), each with a written soundness argument tying its constraints to the reference specification, so the precompile layer preserves the scheme's unforgeability under the assumption that the zkVM's cross-table lookup argument is sound. The criterion then separates the three: the Griffin chip pays against emulation and recovers feasibility (\S5), yet loses to a field-matched control; the field-multiplication chip does not pay, since the host's $256$-bit modular-multiplication precompile already serves its arithmetic; and the power-residue symbol, the one irreducibly-large-prime primitive, pays $10.33\times$ on total committed trace area.
